\section{Parity-pattern guide to the factorial constructions}\label{app:parity-summary}

For quick reference, the table below lists every odd/even pattern through six
inputs and points to a proved construction for the factorial target.  The
``additional block'' column records parity-specific strengthenings that are
explicitly used or singled out in the manuscript.  Universal theorems, such
as the further four-input guarantees of \Cref{app:four-refinements}, apply to
every parity row at that input size and are therefore not repeated in every
cell.

\small
\begin{longtable}{@{}ccl>{\raggedright\arraybackslash}p{0.31\textwidth}>{\raggedright\arraybackslash}p{0.28\textwidth}@{}}
\toprule
$n$ & odd & parity pattern & factorial construction & additional displayed block(s) / proof source \\
\midrule
\endfirsthead
\toprule
$n$ & odd & parity pattern & factorial construction & additional displayed block(s) / proof source \\
\midrule
\endhead
1&0&$0\OO1\EE$&multiple of $1$&the even singleton also gives $2$; \Cref{app:small}\\
1&1&$1\OO0\EE$&multiple of $1$&\Cref{app:small}\\
\midrule
2&0&$0\OO2\EE$&multiple of $2$&also $6,8$; \Cref{app:small}\\
2&1&$1\OO1\EE$&multiple of $2$&universal $3$ also holds; \Cref{app:small}\\
2&2&$2\OO0\EE$&multiple of $2$&also $3,4$; \Cref{app:small}\\
\midrule
3&0&$0\OO3\EE$&$12$ (hence $3!$)&also $14,18,48$; \Cref{app:three12,app:three18,app:three48}\\
3&1&$1\OO2\EE$&$12$ (hence $3!$)&\Cref{app:three12}\\
3&2&$2\OO1\EE$&$12$ (hence $3!$)&\Cref{app:three12}\\
3&3&$3\OO0\EE$&$12$ (hence $3!$)&also $14,18$; \Cref{app:three12,app:three18}\\
\midrule
4&0&$0\OO4\EE$&$24=4!$&also $48,72$; \Cref{thm:fourblocks,app:four48,app:four72}\\
4&1&$1\OO3\EE$&$24=4!$&also $40,48$; \Cref{thm:fourblocks,app:oneodd40,app:four48}\\
4&2&$2\OO2\EE$&$24=4!$&also $40,48$; \Cref{thm:fourblocks,app:twoodd40,app:four48}\\
4&3&$3\OO1\EE$&$24=4!$&also $60$; \Cref{thm:fourblocks,thm:3o1e60}\\
4&4&$4\OO0\EE$&$24=4!$&also $48,60$; \Cref{thm:fourblocks,app:four48,app:fourodd60}\\
\midrule
5&0&$0\OO5\EE$&$120=5!$&also $480$; \Cref{thm:five120,app:fiveeven}\\
5&1&$1\OO4\EE$&$120=5!$&\Cref{thm:five120}\\
5&2&$2\OO3\EE$&$120=5!$&\Cref{thm:five120}\\
5&3&$3\OO2\EE$&$120=5!$&\Cref{thm:five120}\\
5&4&$4\OO1\EE$&$120=5!$&\Cref{thm:five120}\\
5&5&$5\OO0\EE$&$120=5!$&also $180,240$; \Cref{thm:five120,app:fiveodd}\\
\midrule
6&0&$0\OO6\EE$&$18\cdot10\cdot4=720$&\Cref{app:six0}\\
6&1&$1\OO5\EE$&$20\cdot36$ or $10\cdot2\cdot36$&\Cref{app:six1}\\
6&2&$2\OO4\EE$&$720$ via the exhaustive branches in the proof&\Cref{app:six2}\\
6&3&$3\OO3\EE$&$6\cdot120$ or $12\cdot60$&\Cref{app:six3}\\
6&4&$4\OO2\EE$&$720$ via the exhaustive branches in the proof&\Cref{app:six4}\\
6&5&$5\OO1\EE$&$60\cdot12$ or $20\cdot36$&\Cref{app:six5}\\
6&6&$6\OO0\EE$&$18\cdot10\cdot4=720$&\Cref{app:six6}\\
\bottomrule
\end{longtable}
\normalsize

For $n=5$, the single universal congruence proof of \Cref{thm:five120}
covers all six parity rows at once.  For $n=6$, the seven rows correspond
exactly to the seven exhaustive parity branches proved in \Cref{app:six}.
